% Visual analysis at baseline and follow-up: defines the clean-baseline cohort and the outcome, and
% shows where plaque developed (CGPS_data/plaque_location.md, clean_baseline_plaque_location.py).
\section{Visual Analysis at Baseline and Follow-Up}
\label{sec:data-plaque-location}

Every baseline and follow-up CCTA was analysed visually in Medis.
For each of the 17 coronary segments, the reader recorded the stenosis grade (none, minimal, mild, moderate, severe, or occlusion), the plaque type (calcified, mostly calcified, mostly non-calcified, or non-calcified), high-risk plaque features, and evaluability.
The reads were summarised per participant in the segment involvement score (SIS) and the segment stenosis score (SSS), counting the segments with plaque and summing the stenosis grades, respectively.

All SCORE2 risk factors \cite{score2_2021} were also recorded at both visits: age, sex, current smoking, systolic blood pressure, and total and high-density lipoprotein (HDL) cholesterol, from which non-HDL cholesterol is derived.

A participant had a clean baseline when both SIS and SSS were zero at baseline, that is, when no plaque or stenosis was recorded in any segment.
This held for 1,352 participants, 1,057 with a grade per segment and 295 with the summary scores only.
Of these, [n] were excluded due to insufficient image quality, [criterion], leaving [n] participants.
\clarify{image-quality criterion (both from not a good enough scan at baseline or follow up. I will compute these numbers again after the further restriction)}
The restrictions are summarised in Figure~\ref{fig:data-flowchart}.

% Cohort flow chart. Needs \usepackage{tikz} and \usetikzlibrary{positioning, arrows.meta, calc} in the preamble.
% [n] are filled in once the image-quality restriction is computed; per-artery events change after it.
\begin{figure}[htbp]
  \centering
  \small
  \begin{tikzpicture}[
      node distance=6mm and 8mm,
      base/.style={rounded corners=3pt, align=center, inner sep=5pt, minimum height=12mm, very thick},
      box/.style={base, draw=blue!60!black, fill=blue!10, text width=30mm},
      excl/.style={base, draw=red!65!black, fill=red!8, text width=30mm, minimum height=9mm},
      art/.style={base, text width=20mm, minimum height=9mm},
      arr/.style={-{Stealth[length=2.2mm]}, thick, gray!70!black},
    ]
    \node[box] (all) {BL and FU CCTA\\\mbox{$n = 2{,}116$}};
    \node[box, right=of all] (clean) {Clean baseline\\\mbox{$n = 1{,}352$}};
    \node[box, right=of clean] (iq) {Good image quality\\\mbox{$n = $ [n]}};

    \node[art, draw=teal!70!black, fill=teal!12, right=12mm of iq] (lcx) {LCx\\\mbox{[n] / [n]}};
    \node[art, draw=orange!80!black, fill=orange!15, above=4mm of lcx] (lad) {LAD\\\mbox{[n] / [n]}};
    \node[art, draw=violet!70!black, fill=violet!12, below=4mm of lcx] (rca) {RCA\\\mbox{[n] / [n]}};

    \node[excl, below=9mm of $(all)!0.5!(clean)$] (x1) {Plaque at BL\\\mbox{$n = 764$}};
    \node[excl, below=9mm of $(clean)!0.5!(iq)$] (x2) {Poor image quality\\\mbox{$n = $ [n]}};

    \draw[arr] (all) -- (clean);
    \draw[arr] (clean) -- (iq);
    \draw[arr, red!65!black] ($(all.east)!0.5!(clean.west)$) -- (x1);
    \draw[arr, red!65!black] ($(clean.east)!0.5!(iq.west)$) -- (x2);

    \coordinate[right=6mm of iq] (split);
    \draw[thick, gray!70!black] (iq) -- (split);
    \draw[arr] (split) |- (lad);
    \draw[arr] (split) -- (lcx);
    \draw[arr] (split) |- (rca);
  \end{tikzpicture}
  \caption{Restriction of the Herlev--{\O}sterbro data to the study cohort. A clean baseline means that no plaque or stenosis was recorded in any coronary segment at baseline. The same cohort entered all three artery models, shown as participants / participants with new plaque (minimal stenosis or worse) anywhere in the artery's proximal, mid, or distal segment at follow-up. BL, baseline; FU, follow-up.}
  \label{fig:data-flowchart}
\end{figure}

At follow-up, 668 of the 1,352 participants (49.4\%) had developed plaque in at least one segment, with plaque defined as minimal stenosis or worse.
Of these, 626 (93.7\%) had plaque in at least one of the three main vessels (LAD, LCx, and RCA), whereas only 42 (6.3\%) developed it solely in the LM or the side branches.
Table~\ref{tab:data-plaque-location} gives the distribution over vessels, with the 17 segments grouped into 11 vessels and a vessel counted once per participant regardless of how many of its segments (proximal, mid, or distal) had plaque.

\begin{table}[htbp]
  \centering
  \caption{New plaque at follow-up per vessel among the 1,352 participants with a clean baseline. The 17 segments are grouped into 11 vessels, each counted once per participant, and 1,552 vessels had plaque at follow-up. Participant percentages are relative to the 1,352 participants and vessel percentages to the 1,552 vessels with plaque. Other comprises the LM and the side branches. A participant can count in several vessels, and percentages may not sum to 100 due to rounding.}
  \label{tab:data-plaque-location}
  \small
  \begin{tabular}{llrr}
    \toprule
    Vessel & Segments & Participants, $n$ (\%) & Vessels with plaque, $n$ (\%) \\
    \midrule
    LAD & pLAD, mLAD, dLAD & 544 (40.2) & 544 (35.1) \\
    LCx & pLCx, mLCx, dLCx & 208 (15.4) & 208 (13.4) \\
    RCA & pRCA, mRCA, dRCA & 279 (20.6) & 279 (18.0) \\
    Other & LM, side branches & [n] ([\%]) & 521 (33.6) \\
    \midrule
    Total & & 668 (49.4) & 1,552 (100) \\
    \bottomrule
  \end{tabular}
\end{table}
\clarify{participants with plaque in the LM or any side branch (union, not 167 + 239, since a participant can have both); not in \texttt{plaque\_location.md}.}
